\item Reconoce las relaciones e interacciones entre los elementos y los factores del clima como base para comprender, de manera general, la distribución de las regiones naturales en la Tierra y analizar la biodiversidad en el mundo.
\item Indaga y analiza de manera crítica los cambios ocurridos en el clima, sus causas y consecuencias en México y el mundo.
\item Distingue la distribución de las regiones bioculturales principales en México y el mundo.
\item Identifica que los desastres pueden ser originados por procesos naturales o por las actividades humanas. 
\item Relaciona los efectos ambientales, sociales y económicos de los desastres recientes en México y el mundo, tomando en cuenta la vulnerabilidad de la población.
\item Valora la importancia de consolidar una cultura de prevención de desastres con la participación de instituciones, organismos y sociedad. 
\item Analiza las implicaciones sociales, ambientales y económicas del crecimiento, distribución y composición de la población en diferentes países, con base en información estadística y cartográfica. 
\item Emplea las nociones de concentración y dispersión de la población para explicar los rasgos y problemas del espacio urbano y el rural.
\item Distingue la movilidad como un derecho humano, los tipos de migración y los principales flujos migratorios para comprender los efectos socioeconómicos y culturales en los lugares de atracción y expulsión de población.
\item Distingue los tipos de migración y los principales flujos migratorios para comprender los efectos socioeconómicos y culturales en los lugares de atracción y expulsión de la población.
\item Valora la importancia del espacio en la conformación de las identidades juveniles.
\item Debate acerca de la multicausalidad de los conflictos territoriales en México y el mundo, la importancia de la ubicación geográfica de las partes involucradas y las consecuencias ambientales, sociales, económicas y políticas en México y el mundo.
\item Promueve alternativas de resolución justas y pacíficas a los conflictos territoriales que afectan a la comunidad, a México y al mundo.
\item Analiza y relaciona distintos procesos productivos sustantivos en la conformación social, económica y espacial de las sociedades a nivel mundial, para identificar sus contradicciones y desigualdades.
\item Compara la producción agrícola, ganadera, pesquera y forestal en diferentes regiones del mundo.